\clearpage
\section{真人纸条与会话建立}
\label{sec:design-real-chat}
真人交流以短纸条作为低压力的首次邀请，由真实收件人决定是否开始聊天。
投递要求对方为已收藏且仍存在的其他账号、允许接收纸条，双方尚无会话且任一方向都没有待回应纸条。
纸条与回信均须为非空的100字以内文本；每账号每天最多向真人投递2张。
接收开关控制新纸条投递，已有会话继续交流。

\begin{figure}[H]
\centering
\begin{tikzpicture}
\node[reportbox,text width=66mm] (a) at (-1.6,0) {收件人查看仍待回应的纸条};
\node[reportbox,text width=66mm] (b) at (-1.6,-1.15) {接受还是拒绝？};
\node[reportbox,text width=66mm] (c) at (-1.6,-2.3) {接受：填写有效回信\\确认开始交流};
\node[reportbox,text width=66mm] (d) at (-1.6,-3.8) {接受纸条并建立会话\\原纸条与回信成为首轮消息};
\node[reportentity,text width=27mm] (r) at (4.9,-1.15) {拒绝邀请\\不建立会话};
\node[reportentity,text width=27mm] (err) at (4.9,-3.8) {建立失败\\保留待回应状态};
\node[reportentity,text width=66mm] (e) at (-1.6,-5.25) {更新纸条与会话展示};
\draw[reportarrow] (a.south) -- (b.north);
\draw[reportarrow] (b.south) -- (c.north);
\draw[reportarrow] (b.east) -- (r.west);
\draw[reportarrow] (c.south) -- (d.north);
\draw[reportarrow] (d.south) -- (e.north);
\draw[reportarrow] (d.east) -- (err.west);
\draw[reportarrow] (r.east) -- (6.6,-1.15) -- (6.6,-6.2) -- (-1.6,-6.2) -- (e.south);
\end{tikzpicture}
\caption{真人回应纸条与开始交流的选择流程}\label{fig:design-real-flow}
\end{figure}

\subsection{回应选择与会话准入}
收件人回应待处理纸条，拒绝后结束本次邀请；接受时将纸条状态、会话和首轮消息整体保存。
原纸条与回信接入聊天，使首次邀请自然延续为对话。

\subsection{持续交流与页面同步}
建立会话后，双方参与者可以查看消息，并发送非空的500字以内文本。
聊天页按交流顺序展示近期消息，定期同步对方的新内容，并保留正在编辑的草稿与输入焦点。
例如对方在用户输入时回复，页面补充新消息而不打断表达。

